\documentclass[10pt,a4paper]{amsart}
\usepackage{amsmath,amssymb,amsthm,graphicx}
\usepackage{verbatim}
\newenvironment{reference}{\comment}{\endcomment}
%\newenvironment{reference}{}{}
\newenvironment{slogan}{\comment}{\endcomment}
\newenvironment{history}{\comment}{\endcomment}

% For commutative diagrams we use Xy-pic
\usepackage[all]{xy}

% We use 2cell for 2-commutative diagrams.
\xyoption{2cell}
% Theorem environments.
%
\theoremstyle{plain}
\newtheorem{theorem}[subsection]{Theorem}
\newtheorem{proposition}[subsection]{Proposition}
\newtheorem{lemma}[subsection]{Lemma}

\theoremstyle{definition}
\newtheorem{definition}[subsection]{Definition}
\newtheorem{example}[subsection]{Example}
\newtheorem{exercise}[subsection]{Exercise}
\newtheorem{situation}[subsection]{Situation}

\theoremstyle{remark}
\newtheorem{remark}[subsection]{Remark}
\newtheorem{remarks}[subsection]{Remarks}

\numberwithin{equation}{subsection}

% Macros
%
\def\lim{\mathop{\mathrm{lim}}\nolimits}
\def\colim{\mathop{\mathrm{colim}}\nolimits}
\def\Spec{\mathop{\mathrm{Spec}}}
\def\Hom{\mathop{\mathrm{Hom}}\nolimits}
\def\Ext{\mathop{\mathrm{Ext}}\nolimits}
\def\SheafHom{\mathop{\mathcal{H}\!\mathit{om}}\nolimits}
\def\SheafExt{\mathop{\mathcal{E}\!\mathit{xt}}\nolimits}
\def\Sch{\mathit{Sch}}
\def\Mor{\mathop{\mathrm{Mor}}\nolimits}
\def\Ob{\mathop{\mathrm{Ob}}\nolimits}
\def\Sh{\mathop{\mathit{Sh}}\nolimits}
\def\NL{\mathop{N\!L}\nolimits}
\def\CH{\mathop{\mathrm{CH}}\nolimits}
\def\proetale{{pro\text{-}\acute{e}tale}}
\def\etale{{\acute{e}tale}}
\def\QCoh{\mathit{QCoh}}
\def\Ker{\mathop{\mathrm{Ker}}}
\def\Im{\mathop{\mathrm{Im}}}
\def\Coker{\mathop{\mathrm{Coker}}}
\def\Coim{\mathop{\mathrm{Coim}}}

% Boxtimes
%
\DeclareMathSymbol{\boxtimes}{\mathbin}{AMSa}{"02}

%
% Macros for moduli stacks/spaces
%
\def\QCohstack{\mathcal{QC}\!\mathit{oh}}
\def\Cohstack{\mathcal{C}\!\mathit{oh}}
\def\Spacesstack{\mathcal{S}\!\mathit{paces}}
\def\Quotfunctor{\mathrm{Quot}}
\def\Hilbfunctor{\mathrm{Hilb}}
\def\Curvesstack{\mathcal{C}\!\mathit{urves}}
\def\Polarizedstack{\mathcal{P}\!\mathit{olarized}}
\def\Complexesstack{\mathcal{C}\!\mathit{omplexes}}
% \Pic is the operator that assigns to X its picard group, usage \Pic(X)
% \Picardstack_{X/B} denotes the Picard stack of X over B
% \Picardfunctor_{X/B} denotes the Picard functor of X over B
\def\Pic{\mathop{\mathrm{Pic}}\nolimits}
\def\Picardstack{\mathcal{P}\!\mathit{ic}}
\def\Picardfunctor{\mathrm{Pic}}
\def\Deformationcategory{\mathcal{D}\!\mathit{ef}}
\usepackage{cleveref}
\crefname{originallemma}{original lemma}{original lemmas}
\Crefname{originallemma}{Original lemma}{Original lemmas}
\makeatletter
\newlabel{lemma-play-with-topologies}{{00Z7}{source}{}{sourcepointer-1}{}}
\newlabel{lemma-play-with-topologies@cref}{{[originallemma][0][]00Z7}{[1][0][]source}{}{}{}}
\makeatother
\title{The finest topology making a given family of presheaves into sheaves}
\author{The Stacks Project Authors}
\date{stacks-project:sites.tex:local-snapshot-20261012; GNU FDL 1.2 or later}
\begin{document}
\maketitle
\noindent\textit{Editorial note.} All delivered mathematical text is the author's own, copied from the pinned Stacks Project chapter file sites-cohomology.tex; no statement or proof step was rewritten, reordered or removed.\par\medskip
\noindent\textit{Editorial note.} The proof is the author's complete argument for the two assertions: the reduction to a single presheaf, the axioms of a topology and the statement about the finest topology are all written out. The closing sentence 'We leave it to the reader to verify that the construction is inverse to the natural map displayed above' concerns a routine verification and is preserved verbatim inside the counted span.\par\medskip
\noindent\textit{Editorial note.} Every labelled reference occurring inside the counted spans is carried in assembly.external\_references as a non-load-bearing navigation pointer (same-chapter ones at their definition site, cross-chapter ones at the invocation inside the counted span), so no load-bearing mathematics is replaced by a stub. No bibliography citation occurs inside the counted spans.\par\medskip
\noindent\textit{Editorial note.} The checked chapter list for same-name collisions is the set of stacks chapters staged in this repository: chow, cohomology, divisors, etale-cohomology, flat, more-algebra, more-morphisms, morphisms, spaces-more-morphisms, varieties together with this card's chapter sites-cohomology.\par\medskip
\noindent\textit{Editorial note.} The preamble is the chapter's own retained minimal preamble (sources/stacks-sites-cohomology/retained-preamble.tex), a verbatim excerpt of the native preamble: the Xy-pic setup and the native theorem declarations with the mathematical macros.\par\medskip
\noindent\textit{Allowed standard prerequisites (editorial).}\begin{itemize}
\item Intersections of topologies. Same chapter, Lemma lemma-play-with-topologies. Shows that an intersection of topologies is a topology, which reduces the statement to a single presheaf.
\item The axioms of a topology. Same chapter, Definition definition-topology. Supplies the axioms verified in the proof.
\end{itemize}\medskip
\ifdefined\hypertarget\hypertarget{sourcepointer-1}{}\fi
\noindent\textit{Source reference.} Original lemma 00Z7: sources/stacks-sites/pinned-original-source.tex, source lines 11362--11362. Navigation pointer invoked at proof line 11362; the counted span contains the pointer only and the cited result is not reproduced inside it, so the pointer is registered as a non-load-bearing navigation stub in the established form for Stacks material. No bibliography citation is involved.\par\medskip
\begin{lemma}
\label{lemma-topology-presheaves-sheaves}
Let $\mathcal{C}$ be a category. Let $\{ \mathcal{F}_i \}_{i\in I}$ be a
collection of presheaves of sets on $\mathcal{C}$. For each
$U \in \Ob(\mathcal{C})$ denote
$J(U)$ the set of sieves $S$ on $U$ with the following property:
For every morphism $V \to U$, the maps
$$
\Mor_{\textit{PSh}(\mathcal{C})}(h_V, \mathcal{F}_i)
\longrightarrow
\Mor_{\textit{PSh}(\mathcal{C})}(S \times_U V, \mathcal{F}_i)
$$
are bijective for all $i \in I$. Then $J$ defines a
topology on $\mathcal{C}$. This topology is the finest
topology in which all of the $\mathcal{F}_i$ are sheaves.
\end{lemma}

\begin{proof}
First, we note that if $J$ is a topology, then the last statement of
the lemma immediately follows.

\medskip\noindent
For each $i \in I$ and $U \in \Ob(\mathcal{C})$ let $J_i(U)$
be the set of sieves $S$ on $U$ with the following property:
For every morphism $V \to U$, the map
$$
\Mor_{\textit{PSh}(\mathcal{C})}(h_V, \mathcal{F}_i)
\longrightarrow
\Mor_{\textit{PSh}(\mathcal{C})}(S \times_U V, \mathcal{F}_i)
$$
is bijective. Note that $J(U) = \bigcap J_i(U)$. Thus if we prove
that $J_i$ defines a topology, then we conclude by
Lemma \ref{lemma-play-with-topologies}.
Thus we reduce to the case discussed in the next paragraph.

\medskip\noindent
Proof that $J$ is a topology when our family is given by a single
presheaf $\mathcal{F}$. The first and third axioms of
a topology are immediately verified. Thus, assume that
we have an object $U$, and sieves $S, S'$ of $U$
such that $S \in J(U)$, and for all $V \to U$ in $S(V)$
we have $S' \times_U V \in J(V)$. We have to show that
$S' \in J(U)$. In other words, we have to show that for
any $f : W \to U$, the map
$$
\mathcal{F}(W) =
\Mor_{\textit{PSh}(\mathcal{C})}(h_W, \mathcal{F})
\longrightarrow
\Mor_{\textit{PSh}(\mathcal{C})}(S' \times_U W, \mathcal{F})
$$
is bijective. Pick an element $\varphi \in
\Mor_{\textit{PSh}(\mathcal{C})}(S' \times_U W, \mathcal{F})$.
We will construct a section $s \in \mathcal{F}(W)$
mapping to $\varphi$.

\medskip\noindent
Suppose $\alpha : V \to W$ is an element of $S \times_U W$.
According to the definition of pullbacks we see that
the composition $f \circ \alpha : V \to W  \to U$ is in $S$. Hence
$S' \times_U V$ is in $J(V)$ by assumption on the pair
of sieves $S, S'$. Now we have a commutative diagram
of presheaves
$$
\xymatrix{
S' \times_U V \ar[r] \ar[d] & h_V \ar[d] \\
S' \times_U W \ar[r] & h_W
}
$$
The restriction of $\varphi$ to $S' \times_U V$
corresponds to an element $s_{V, \alpha} \in \mathcal{F}(V)$.
This we see from the definition of $J$,
and because $S' \times_U V$ is in $J(V)$.
We leave it to the reader to check
that the rule $(V, \alpha) \mapsto s_{V, \alpha}$ defines
an element
$\psi \in
\Mor_{\textit{PSh}(\mathcal{C})}(S \times_U W, \mathcal{F})$.
Since $S \in J(U)$ we see immediately from the definition of $J$
that $\psi$ corresponds to an element $s$ of $\mathcal{F}(W)$.

\medskip\noindent
We leave it to the reader to verify that the construction
$\varphi \mapsto s$ is inverse to the natural map displayed above.
\end{proof}

\end{document}
